\chapter*{Abstract}
\addcontentsline{toc}{chapter}{Abstract}

An automated vehicle approaching a signalized intersection does not need the color of every traffic
light in view; it needs to know whether its intended maneuver---straight, left or right---is
permitted. This thesis studies this \emph{maneuver-level} formulation of traffic signal understanding,
introduced by Zou et al., in which one front-camera image is mapped to a signal color and a stop/go
decision for each maneuver.

The work has three objectives. First, since the original benchmark is not public, it rebuilds the
maneuver-level labels from the lane and traffic-element topology of OpenLane-V2 and validates them
against the published label distribution. Second, it builds a structured model, in which a real-time
detector (RF-DETR) classifies each light by the maneuvers it controls and an explicit rule assigns colors
to maneuvers. Third, it tests transfer across countries, traffic rules and cameras, using new compatible
labels for the DriveU (Germany) and LISA (United States) traffic light datasets and a per-light
evaluation on ATLAS.

The rebuilt labels reproduce the published distribution for almost every class and indicate that the
original benchmark combines both OpenLane-V2 subsets under US rules. The structured model reaches a
Traffic Signal Score (TSS) of 0.899 on subset A and 0.915 on subset B, in the range published for the
end-to-end model. Trained on DTLD it reaches a test TSS of 0.958 at 22\,ms per frame; it transfers well
to LISA by day (0.952) but drops at night (0.863), and its False-Go rate doubles on Argoverse~2 images.
Right-turn signals, yellow lights, small distant lights and country-specific label conventions emerge
as the main open problems, and the safety-relevant errors concentrate exactly where aggregate metrics
do not look.

\vspace{0.5cm}
\noindent\textbf{Keywords:} traffic light recognition, autonomous driving, signal intent, DETR,
RF-DETR, OpenLane-V2, domain shift, traffic rules.
